\subsection{Divisorial smoothing}
\label{sec:divisorial-adjoint-obstruction}
\index{Adjoint measure!divisorial smoothing}
\index{Monge--Ampere mass@Monge--Amp\`ere mass!concentration on a divisor}

For a singular reference, the supporting probability can be unique while smoothings of the solution have different limiting adjoint probabilities. The reason is that the non-pluripolar product discards the mass on a divisor that the smoothings keep. On a product of two projective lines, explicit formulas exhibit both the lost mass and its dependence on the smoothing.

In this subsection, $X=\mathbb P^1\times\mathbb P^1$. Let $\omega_1$ and $\omega_2$ be the pullbacks of the Fubini--Study forms of the two factors, normalized to have integral one on their factors, and put
\[
 \theta=2\omega_1+2\omega_2,\qquad
 \beta=\omega_1+2\omega_2,\qquad
 \omega=\omega_1+\omega_2,\qquad
 \mathrm dV=\omega_1\wedge\omega_2.
\]
Thus $\mathrm dV$ is a probability volume form, $\theta^2=8\mathrm dV$ and $\beta^2=4\mathrm dV$. Let $z$ be the affine coordinate on the first factor, and put
\[
 r=\frac{|z|^2}{1+|z|^2},\qquad \psi=\log r,\qquad
 D=\{0\}\times\mathbb P^1.
\]
The function $r$ extends smoothly to the first factor, with value one at infinity, and $\psi$ is quasi-plurisubharmonic with value $-\infty$ exactly on $D$. The Poincar\'e--Lelong formula gives $\ddc\psi=[D]-\omega_1$, so $\theta+\ddc\psi=\beta+[D]$.

For a smooth positive $(1,1)$-form $\Omega$ on the surface $X$, the linearization $v\mapsto2\ddc v\wedge\Omega/\Omega^2$ of the logarithmic Monge--Amp\`ere operator has the adjoint probability $\rho\Omega^2$, where the Gauduchon multiplier $\rho>0$ satisfies $\ddc(\rho\Omega)=0$ and $\int_X\rho\Omega^2=1$ (\cref{prop:adjoint-gauduchon}). We call $Q=\rho\Omega$ the \emph{adjoint cofactor} of $\Omega$: it is $\ddc$-closed, and $\Omega\wedge Q$ is the adjoint probability.

\Needspace{10\baselineskip}
\begin{lemma}\label{lem:divisorial-adjoint-model}
Write $x=\log|z|^2$ and $\chi(x)=\log(1+\mathrm e^x)$. The envelope $\varphi=P_\theta[\psi]$ is
\begin{equation}\label{eq:divisorial-adjoint-model}
 \varphi=W(x)-2\chi(x),\qquad
 W(x)=\begin{cases}x+2\log2,&x\leq0,\\2\chi(x),&x\geq0.\end{cases}
\end{equation}
It satisfies $\psi\leq\varphi\leq0$, $0\leq\varphi-\psi\leq2\log2$, and the model condition \eqref{eq:als-model-condition} with $\varepsilon=3/4$.
\end{lemma}
\begin{proof}
The function $W$ is convex: both one-sided derivatives at zero equal one, and the derivative of $2\chi$ increases from one to two on $[0,\infty)$. Since $\omega_1=\ddc\chi(x)$ on the affine chart, $W-2\chi$ is $2\omega_1$-plurisubharmonic there. Since $W$ has slope one for $x\leq0$ and agrees with $2\chi$ for $x\geq0$, the function $W-2\chi$ extends to a global $2\omega_1$-plurisubharmonic function, with a logarithmic pole along $D$ and equal to zero where $|z|\geq1$. The asserted bounds follow from $\psi=x-\chi$: the difference $\varphi-\psi$ equals $2\log2-\chi$ for $x\leq0$ and $\chi-x$ for $x\geq0$.

For $C>0$ write $v_C=P_\theta(\min\{\psi+C,0\})$. Rotations of the first factor and the transitive unitary action on the second preserve the whole competitor class, so the envelope is radial and independent of the second factor. The function $V_C=2\chi+v_C$ is then convex in $x$ and satisfies $V_C\leq\chi+x+C$. This upper bound as $x\to-\infty$ forces the slope of $V_C$ at $-\infty$ to be at least one, and by convexity all its slopes are at least one. For $x<0$,
\[
 V_C(x)\leq V_C(0)+x\leq2\log2+x=W(x).
\]
For $x\geq0$, the zero obstacle gives $V_C\leq2\chi=W$. Consequently $v_C\leq\varphi$. If $C\geq2\log2$, the bounds already proved make $\varphi$ a competitor, so $v_C=\varphi$. Letting $C\to\infty$ proves $P_\theta[\psi]=\varphi$. Since $\varphi\sim\psi$ and the ordinary envelope depends only on the singularity type, also $P_\theta[\varphi]=\varphi$.

Finally
\[
 \tfrac34\theta+\ddc\psi
 =\tfrac12\omega_1+\tfrac32\omega_2+[D]\geq0.
\]
Since $\psi\leq\varphi$, the function $\psi$ is a competitor for $P_{(3/4)\theta}(\varphi)$, which is therefore not identically $-\infty$.
\end{proof}

\Needspace{10\baselineskip}
\begin{proposition}\label{prop:divisorial-residual-adjoint}
Every $v\in\PSH(X,\theta)$ with $v\sim\varphi$ is uniquely of the form $v=\psi+w$, where $w\in\PSH(X,\beta)$ is bounded. Its non-pluripolar measure is
\begin{equation}\label{eq:divisorial-residual-mass}
 \mathrm{MA}_\theta(v)=(\beta+\ddc w)^2,\qquad
 \int_X\mathrm{MA}_\theta(v)=4.
\end{equation}
For every real continuous $F$,
\[
 B_\varphi(F)=\log4-\log\int_X\mathrm e^F\mathrm dV,
 \qquad
 \mathcal A_\varphi(F)=
 \left\{\frac{\mathrm e^F\mathrm dV}{\int_X\mathrm e^F\mathrm dV}\right\}.
\]
At $v=\psi$, the non-pluripolar part of $\theta+\ddc\psi$ is $\beta$, whose adjoint cofactor is $Q=\beta/4$, with adjoint probability $\beta\wedge Q=\mathrm dV$.
\end{proposition}
\begin{proof}
On $X\setminus D$, the function $w=v-\psi$ is bounded, because $v\sim\varphi\sim\psi$, and satisfies $\beta+\ddc w\geq0$. After adding local potentials of the closed form $\beta$, \cref{lem:qpsh-removable-pluripolar} extends $w$ across $D$; the extensions agree on overlaps and give a bounded $\beta$-plurisubharmonic function on $X$. Conversely, for every bounded $w\in\PSH(X,\beta)$, the function $\psi+w$ is $\theta$-plurisubharmonic with current $\beta+\ddc w+[D]$, and it has the ordinary type of $\varphi$. The decomposition is unique, since two such functions $w$ agree off $D$ and hence everywhere.

Put $T=\beta+\ddc w$. The bounded-product construction gives
\[
 T^2-\beta^2=\ddc\bigl(w(T+\beta)\bigr),
\]
and Stokes' theorem against the constant function yields $\int_XT^2=4$. The product gives no mass to $D$ by \cref{prop:bounded-products}; we also give a direct cutoff estimate, which will be used again below. Shift $w$ so that $0\leq w\leq M$, and set $f_t=\max\{\psi/t,-1\}$ for $t>0$. Choose smooth convex functions $\sigma_\epsilon$ with derivative in $[0,1]$, equal to $-1$ below $-1-\epsilon$ and to the identity above $-1+\epsilon$, converging uniformly to $\max\{\cdot,-1\}$. Then $f_{t,\epsilon}=\sigma_\epsilon(\psi/t)$ is globally smooth, constant near $D$, and
\[
 f_{t,\epsilon}\longrightarrow f_t\ \text{uniformly},\qquad
 \ddc f_{t,\epsilon}\geq-\omega_1/t.
\]
Testing the distributional identity against $f_{t,\epsilon}$, using positivity of $T+\beta$, and letting $\epsilon\to0$, we obtain
\begin{equation}\label{eq:divisorial-cutoff-bound}
 \int_X(-f_t)T^2
 \leq\int_X(-f_t)\beta^2
   +\frac Mt\int_X\omega_1\wedge(T+\beta)
 =\int_X(-f_t)\beta^2+\frac{4M}{t}.
\end{equation}
where the last integral equals four by Stokes' theorem. The right-hand side tends to zero as $t\to\infty$, since $f_t\to0$ off $D$ and $\beta^2(D)=0$; as $-f_t=1$ on $D$, this gives $T^2(D)=0$. Off $D$ the currents $\theta+\ddc v$ and $T$ agree, so locality and the fact that the extended measure puts no mass on the pole set prove \eqref{eq:divisorial-residual-mass}.

For the Hermitian background $\beta$, \cref{thm:kn-continuous} with density $\mathrm e^F$ gives a bounded solution $w$ and a constant $c>0$, and the correspondence above turns it into a solution $\psi+w$ in the ordinary type of $\varphi$. Every potential of that type has mass four, so integration forces $c=4/\int_X\mathrm e^F\mathrm dV$, which is the formula for $B_\varphi$. Jensen's inequality for the probability $\nu_F=\mathrm e^F\mathrm dV/\int_X\mathrm e^F\mathrm dV$ gives
\[
 B_\varphi(G)\leq B_\varphi(F)-\int_X(G-F)\,\mathrm d\nu_F.
\]
so $\nu_F$ is a supporting probability. For any supporting probability, testing its inequality at $G=F\pm th$ and differentiating at $t=0$ shows that its integral against every $h\in C(X)$ equals that of $\nu_F$.

For a smooth real $h$ and sufficiently small $|t|$, we have $\beta+t\ddc h>0$, so $\psi+th$ stays in the same ordinary type and its residual measure is $(\beta+t\ddc h)^2$, with derivative $2\ddc h\wedge\beta$ at $t=0$. Since $\beta$ is closed, the constant weight $\rho=1/4$ annihilates this derivative after integration, and $Q=\rho\beta$ gives the asserted probability. This normalization uses the residual mass: since $\int_X\theta\wedge Q=3/2$, the cofactor $Q$ is not a normalized Gauduchon current for $\theta$ in the sense of \cref{sec:gauduchon-linear-duality}.
\end{proof}

\Needspace{12\baselineskip}
\begin{theorem}\label{thm:divisorial-smoothing-adjoints}
Let $g$ be a smooth real function on the second factor with $\gamma=\ddc g\geq-\omega_2$, regarded as a function on $X$. For $0<\delta\leq1$ put
\[
 u_{\delta,g}=\log(r+\delta\mathrm e^g),\qquad
 \Omega_{\delta,g}=\theta+\ddc u_{\delta,g},\qquad
 F_{\delta,g}=\log\frac{\Omega_{\delta,g}^2}{4\mathrm dV}.
\]
The forms satisfy $\Omega_{\delta,g}\geq\omega$ and $\int_X\Omega_{\delta,g}^2=8$. For each fixed $g$, as $\delta\downarrow0$ the smooth potentials $u_{\delta,g}$ decrease pointwise to $\psi$, converge in $L^1$ and in the capacity $\operatorname{Cap}_\omega$, and converge smoothly off $D$. Moreover $F_{\delta,g}\to0$ in $L^p(\mathrm dV)$ for every $1\leq p<\infty$. The adjoint cofactors and adjoint probabilities of the forms $\Omega_{\delta,g}$ are
\[
 Q_{\delta,g}=\Omega_{\delta,g}/8,\qquad
 \nu_{\delta,g}=\Omega_{\delta,g}^2/8,
\]
and they have the weak limits
\begin{equation}\label{eq:divisorial-adjoint-limits}
 \begin{split}
 Q_{\delta,g}&\longrightarrow(\beta+[D])/8,\\
 \nu_{\delta,g}&\longrightarrow
 \tfrac12\mathrm dV+\tfrac12[D]\wedge\omega_2
                  +\tfrac18[D]\wedge\ddc g.
 \end{split}
\end{equation}
The limiting probability has mass $1/2$ on $D$ and depends on $g$, although $\mathcal A_\varphi(0)=\{\mathrm dV\}$. The cofactors do not approach $\beta/4$, and their limit is not $\rho\beta$ for any locally integrable scalar $\rho$.
\end{theorem}
\begin{proof}
Smoothness and monotone convergence follow directly from the defining formula, also on $D$, where $\psi=-\infty$. Off $D$, set
\[
 \lambda=\frac{r}{r+\delta\mathrm e^g},\qquad
 h=\psi-g-\log\delta.
\]
The chain rule gives $\Omega_{\delta,g}=A_\delta+R_\delta$, where
\[
 A_\delta=(2-\lambda)\omega_1+2\omega_2+(1-\lambda)\gamma,
 \qquad R_\delta=\lambda(1-\lambda)\mathrm dh\wedge\dc h.
\]
Thus $A_\delta\geq\omega_1+\omega_2$ and $R_\delta\geq0$, and the inequality $\Omega_{\delta,g}\geq\omega$ extends across $D$ by continuity. Closedness of $\theta$ and Stokes' theorem give the mass eight. Since $\Omega_{\delta,g}$ is closed, the Gauduchon condition $\ddc(\rho\Omega_{\delta,g})=0$ reduces to the elliptic equation $\ddc\rho\wedge\Omega_{\delta,g}=0$ on the connected surface, whose smooth solutions are constant by the maximum principle. Normalization against $\Omega_{\delta,g}^2$ gives $\rho=1/8$ and the stated adjoint cofactors and probabilities.

The rank-one form $R_\delta$ has square zero, hence $\Omega_{\delta,g}^2=A_\delta^2+2A_\delta\wedge R_\delta$. The bounded coefficients of $A_\delta$ converge off $D$ to those of $\beta$, which gives $A_\delta^2\to4\mathrm dV$. Write $N=\mathrm d\psi\wedge\dc\psi$. In $A_\delta\wedge R_\delta$, the part in which $R_\delta$ contributes only derivatives along the first factor (the normal part) is
\[
 \lambda(1-\lambda)N\wedge\bigl(2\omega_2+(1-\lambda)\gamma\bigr).
\]
Terms mixing a derivative along the first factor with one along the second vanish for degree reasons. The remaining tangential term is $\lambda(1-\lambda)(2-\lambda)\omega_1\wedge\mathrm dg\wedge\dc g$, which tends to zero by dominated convergence.

The normal coefficients can be computed in polar coordinates. Put $s=\delta\mathrm e^{g}$ and $t=|z|^2$. Then
\[
 \lambda=\frac{t}{s+(1+s)t},\qquad
 N=\frac{\mathrm i\,\mathrm dz\wedge\mathrm d\bar z}{2\pi t(1+t)^2}.
\]
After angular integration the radial measure of $\lambda(1-\lambda)N$ is
\[
 \frac{s\,\mathrm dt}{(1+t)(s+(1+s)t)^2}
 =\frac{\mathrm dy}{(1+sy)(1+(1+s)y)^2},\qquad t=sy.
\]
It is bounded by $\mathrm dy/(1+y)^2$ and tends to that measure, whose integral is one. With an additional factor $1-\lambda$, the rescaled measure becomes $\mathrm dy/(1+(1+s)y)^3$, which is bounded by $\mathrm dy/(1+y)^3$ and has limiting integral $1/2$. The mass of both measures outside every fixed tube around $D$ tends to zero, uniformly in the point of the second factor because $g$ is bounded. Testing against continuous functions therefore gives
\[
 \lambda(1-\lambda)N\longrightarrow[D],\qquad
 \lambda(1-\lambda)^2N\longrightarrow\tfrac12[D].
\]
Combining the normal and the remaining terms proves
\begin{equation}\label{eq:divisorial-square-limit}
 \Omega_{\delta,g}^2\longrightarrow
 4\mathrm dV+[D]\wedge(4\omega_2+\gamma).
\end{equation}
Dividing by eight gives the second line of \eqref{eq:divisorial-adjoint-limits}. The limit is positive since $4\omega_2+\gamma\geq3\omega_2$, and its mass on $D$ is $4/8=1/2$ by Stokes' theorem on the second factor. Taking $g=0$ and a small nonconstant smooth $g$ gives different limits, since $\ddc g$ vanishes identically only for constant $g$.

The difference of the potentials satisfies
\[
 0\leq u_{\delta,g}-\psi
 \leq\log\bigl(1+\delta\mathrm e^{\max g}/r\bigr)
\]
off $D$. The right-hand side tends to zero and is dominated by an integrable logarithmic pole, which proves $L^1$ convergence. Continuity of distributional differentiation then gives the limit of the cofactors. A locally integrable function times the smooth form $\beta$ has absolutely continuous coefficients, so it cannot produce the singular coefficients of $[D]$.

For convergence in capacity, use the capacity $\operatorname{Cap}_\omega$ of \cref{sec:hermitian-capacity}; after subtracting one, its test potentials range over $v\in\PSH(X,\omega)$ with $-1\leq v\leq0$. For such $v$, apply the cutoff argument of \eqref{eq:divisorial-cutoff-bound} with $\beta$ replaced by $\omega$ and $w=v+1$. Stokes' theorem gives $\int_X\omega_1\wedge(2\omega+\ddc v)=2$, so
\[
 \operatorname{Cap}_\omega\{\psi<-t\}
 \leq\int_X(-f_t)\omega^2+2/t\longrightarrow0.
\]
The same bound shows that $D$ has capacity zero. For $\eta>0$, the set of points off $D$ where $u_{\delta,g}-\psi>\eta$ is contained in
\[
 \left\{r<\frac{\delta\mathrm e^{\max g}}{\mathrm e^\eta-1}\right\}.
\]
a sublevel set of $\psi$ whose level tends to $-\infty$ as $\delta\downarrow0$. The capacity bound for such sublevel sets, together with $\operatorname{Cap}_\omega(D)=0$, proves convergence in capacity.

Finally, $\Omega_{\delta,g}\geq\omega$ gives $F_{\delta,g}\geq-\log2$. In the affine chart $|\mathrm d\psi|_\omega^2\leq C/r$, while $g$ and its first two derivatives are bounded. The decomposition above yields $\Omega_{\delta,g}\leq C_g(1+r^{-1})\omega$, hence
\[
 |F_{\delta,g}|\leq C_g+2|\log r|.
\]
Every finite power of this logarithmic bound is integrable. Smooth convergence off $D$ gives $F_{\delta,g}\to0$ there, and dominated convergence proves convergence in $L^p$.
\end{proof}

The exponential densities satisfy $\int_X\mathrm e^{F_{\delta,g}}\mathrm dV=2$ for every $\delta$, whereas their pointwise limit has integral one; they are therefore not uniformly integrable. The smoothings leave the ordinary type of $\varphi$ and keep the mass eight of bounded potentials instead of the mass four of that type. The example thus shows which concentration a recovery theorem must exclude. The hypotheses of \cref{prop:weak-adjoint-stability} (uniform convergence) and of \cref{thm:normalization-bounded-lone} (uniformly bounded densities converging in $L^1$) are stronger than the convergence in every $L^p$ obtained here. In the fixed divisorial type, a common bound on the residual potentials gives a direct recovery estimate instead.

\Needspace{12\baselineskip}
\begin{corollary}\label{cor:divisorial-residual-recovery}\index{Adjoint measure!fixed-divisor recovery}
Let $w_j,w\in\PSH(X,\beta)$ satisfy $|w_j|,|w|\leq M$ for a fixed $M\geq1$, and suppose that $w_j\to w$ in capacity. Put
\[
 \begin{gathered}
 T_j=\beta+\ddc w_j,\qquad T=\beta+\ddc w,\\
 Q_j=T_j/4,\quad Q=T/4,\qquad
 \nu_j=T_j^2/4,\quad\nu=T^2/4.
 \end{gathered}
\]
Then $Q_j\rightharpoonup Q$ as currents and $\nu_j\rightharpoonup\nu$ as probability measures. More precisely, for real $f\in C^\infty(X)$ choose $C_f\geq0$ with $-C_f\beta\leq\ddc f\leq C_f\beta$. For every $\eta>0$, writing $E_{j,\eta}=\{|w_j-w|>\eta\}$, one has
\begin{equation}\label{eq:divisorial-residual-recovery}
 \left|\int_X f\,(\mathrm d\nu_j-\mathrm d\nu)\right|
 \leq C_f\bigl(2\eta+2M^2\operatorname{Cap}_\beta(E_{j,\eta})\bigr).
\end{equation}
If $T_j$ is smooth and strictly positive, then $Q_j$ and $\nu_j$ are its adjoint cofactor and adjoint probability. If $\psi+w$ solves \eqref{eq:relative-normalization} for some $F\in C(X)$, the limit $\nu$ is the unique member of $\mathcal A_\varphi(F)$.
\end{corollary}
\begin{proof}
The capacity is the one of \cref{sec:hermitian-capacity}, and $\beta\leq2\omega$ gives $\operatorname{Cap}_\beta\leq4\operatorname{Cap}_\omega$. Set $a_j=(w_j+M)/(2M)$ and $S_j=\beta+\ddc a_j$. Then $0\leq a_j\leq1$ and
\[
 \begin{gathered}
 S_j=\left(1-\frac1{2M}\right)\beta+\frac1{2M}T_j,\\
 S_j^2\geq\left(\frac1M-\frac1{2M^2}\right)\beta\wedge T_j
 \geq\frac1{2M}\beta\wedge T_j.
 \end{gathered}
\]
Positivity of bounded mixed products justifies these inequalities. Thus $\int_E\beta\wedge T_j\leq2M\operatorname{Cap}_\beta(E)$ for every Borel set $E$, and the same bound holds for $T$. By closedness and Stokes' theorem, $\int_X\beta\wedge(T_j+T)=8$. Now test the identity
\[
 T_j^2-T^2=\ddc\bigl((w_j-w)(T_j+T)\bigr)
\]
against $f$. On $X\setminus E_{j,\eta}$ the potential difference has absolute value at most $\eta$, and on $E_{j,\eta}$ at most $2M$. Consequently
\[
 \begin{aligned}
 \left|\int_Xf(T_j^2-T^2)\right|
 &\leq C_f\left(8\eta+2M\int_{E_{j,\eta}}\beta\wedge(T_j+T)\right)\\
 &\leq C_f\left(8\eta+8M^2\operatorname{Cap}_\beta(E_{j,\eta})\right).
 \end{aligned}
\]
Dividing by four proves \eqref{eq:divisorial-residual-recovery}. Letting first $j\to\infty$ with $\eta$ fixed and then $\eta\downarrow0$ shows that $\int_Xf\,\mathrm d\nu_j\to\int_Xf\,\mathrm d\nu$. All $\nu_j$ and $\nu$ have mass one by \cref{prop:divisorial-residual-adjoint}, and smooth functions are uniformly dense in $C(X)$, so $\nu_j\rightharpoonup\nu$.

The constant test potential in the definition of capacity gives $4\mathrm dV(E)\leq\operatorname{Cap}_\beta(E)$. Splitting the potential difference over the same sets yields
\[
 \|w_j-w\|_{L^1(\mathrm dV)}
 \leq\eta+\tfrac M2\operatorname{Cap}_\beta(E_{j,\eta}).
\]
Taking the limits in the same order proves $L^1$ convergence, and hence convergence of the currents $Q_j$ by continuity of distributional differentiation. The limiting cofactor satisfies $\ddc Q=0$, $\int_X\beta\wedge Q=1$, and $T\wedge Q=\nu$ as a bounded product. For smooth real $h$, the signed measure $2\ddc h\wedge T$ is the formal linearization of the residual operator at $w$ in the direction $h$, and its integral against the constant $1/4$ vanishes by Stokes' theorem. It is the derivative of the residual measure along $w+th$ whenever these functions remain $\beta$-plurisubharmonic.

If $T_j$ is smooth and strictly positive, the maximum principle on the connected surface makes its Gauduchon multiplier constant, and the mass four normalizes it to $1/4$. Finally, if $\psi+w$ solves the equation at $F$, then $\nu=T^2/4=\mathrm e^F\mathrm dV/\int_X\mathrm e^F\mathrm dV$, which is the unique supporting probability by \cref{prop:divisorial-residual-adjoint}.
\end{proof}

In the corollary, the potentials $\psi+w_j$ keep their pole along $D$, and the adjoint objects refer to the smooth residual forms $T_j$. We do not know whether a similar recovery holds for more general singular references, for unbounded residual sequences, or when the Gauduchon multipliers are not constant.
